\documentclass[10pt,a4paper]{amsart}
\usepackage[margin=19mm]{geometry}
\usepackage{amsmath,amssymb,mathtools}
\usepackage[scr=rsfs,cal=euler]{mathalfa}
\usepackage{xfrac}
\usepackage[unicode,hidelinks,bookmarks=false]{hyperref}
\newcommand{\ie}{i.e.\ }
\newcommand{\cf}{cf.\ }
\newcommand{\resp}{respectively\ }
\newcommand{\Subsection}{\subsection}
\providecommand{\nobreakdash}{\nobreak\hbox{-}}
\setlength{\emergencystretch}{2em}
\multlinegap0pt
\usepackage{bm}
\usepackage{bbm}
\usepackage{dsfont}
\usepackage{xspace}
\usepackage{cancel}
\usepackage{mathtools}
\usepackage{cleveref}
\usepackage{amsthm}
\makeatletter
\@ifundefined{defin}{\newtheorem{defin}{Defin}}{}
\@ifundefined{cor}{\newtheorem{cor}[defin]{Corollary}}{}
\@ifundefined{lem}{\newtheorem{lem}[defin]{Lemma}}{}
\@ifundefined{thm}{\newtheorem{thm}[defin]{Theorem}}{}
\makeatother
\newcommand{\C}{{\mathbb C}}
\newcommand{\N}{{\mathbb N}}
\newcommand{\cF}{{\mathcal F}}
\newcommand{\cH}{{\mathcal H}}
\newcommand{\cR}{{\mathcal R}}
\newcommand{\nl}{\operatorname{nl}}
\newcommand{\ord}{\operatorname{ord}}
\newcommand{\pp}{\mathbb{P}}
\newcommand{\rk}{\operatorname{rk}}
\title[Geometry of Adjoint Hypersurfaces for Polytopes]{Geometry of Adjoint Hypersurfaces for Polytopes}
\author{Clemens Br\"user, Julian Weigert}
\date{Source: arXiv:2511.13537v1 (CC BY 4.0)}
\begin{document}
\maketitle

\noindent\emph{Source and licence.} Geometry of Adjoint Hypersurfaces for Polytopes. Version arXiv:2511.13537v1 (CC BY 4.0), distributed under the Creative Commons Attribution 4.0 International licence, as recorded in the arXiv licence metadata for this exact version and shown on the abstract page https://arxiv.org/abs/2511.13537v1. Full rights notice: licenses/arxiv-2511.13537-CC-BY-4.0.txt.\par\smallskip

\noindent\textit{Editorial scope.} Counted unit: the Thm whose source label is \texttt{thm:pointresidual} in the article (source0.tex lines 1053--1063, source label \texttt{thm:pointresidual}), delivered together with the article's own complete original proof of it (source0.tex lines 1126--1182, one \texttt{proof} environment of 57 lines). No mathematics is written, repaired or re-derived: statement and proof are the article's own text line for line. The proof is detached from the statement in the source: the article states the result at lines 1053--1063 and prints its proof at lines 1126--1182, and the proof environment's own header names the statement, so the association is the article's own. Required context retained uncounted: cor: van\_order\_Mult (lines 1091--1097); lem: convpol\_noninc (lines 1113--1117). Every definition, display and label of those lines is retained unchanged. Omitted from this package and not redistributed: 1--1052, 1064--1090, 1098--1112, 1118--1125, 1183--1311. Nothing inside a retained excerpt is omitted or altered: every retained body is delivered exactly as the article's lines are written, so no omission receipt of any kind accompanies this package. Citations: the retained excerpts carry no citation command at all, so no bibliography entry of the article is transcribed and no reference is redistributed. Reference adaptation: none was needed and none was made; every reference inside the delivered unit resolves inside it, so no unresolved reference or citation is produced. Printed slips kept literal: no printed word, symbol, hypothesis or number of the counted statement or of its proof was altered. Layout and class: the article's own class is \texttt{amsart} with packages including amsfonts, mathrsfs, amsthm, amssymb, tikz-cd, tikz, hyperref, cleveref, url, mathtools, enumitem; the wrapper is the lane's pinned \texttt{amsart} editorial wrapper at 10pt on A4, which restates the theorem-like environments the retained text needs and adds the article's own macro definitions that the retained text needs (\texttt{\textbackslash C}, \texttt{\textbackslash N}, \texttt{\textbackslash cF}, \texttt{\textbackslash cH}, \texttt{\textbackslash cR}, \texttt{\textbackslash nl}, \texttt{\textbackslash ord}, \texttt{\textbackslash pp}, \texttt{\textbackslash rk}), with extra wrapper packages (bm, bbm, dsfont, xspace, cancel, mathtools, cleveref). The delivered numbering is the unit's own. Assets: no retained span contains a figure environment, a graphics-inclusion command, a PDF-page inclusion, a table or a TikZ picture; no image file is copied into this package, the package declares no asset dependency and no image file is redistributed with it. Licence: version arXiv:2511.13537v1 of this article is distributed under the Creative Commons Attribution 4.0 International licence, as recorded in the arXiv licence metadata for this exact version and shown on the abstract page https://arxiv.org/abs/2511.13537v1. A full rights notice is delivered at licenses/arxiv-2511.13537-CC-BY-4.0.txt. Characters: every retained excerpt is pure ASCII, so the delivered engine, XeLaTeX with its default Latin Modern text font, reports no missing character.
\begin{cor} \label{cor: van_order_Mult}
 Fix $D,k\in \N$, let $L\subset \pp^{n}$ be a hyperplane and let $f \in \C[x_0, \dots, x_{n}]_D$ be homogeneous of degree $D$. Furthermore let $H_1, \dots, H_k$ be pairwise distinct hyperplanes in $L$. Assume that there exists $m \in \N$ with
 \begin{equation*}
     \sum_{i=1}^k (\mu_{H_i}(f)-(m-1)) > D - (m -1).
 \end{equation*}
 Then $\mu_L(f) \geq m$.
\end{cor}

\begin{lem}
\label{lem: convpol_noninc}
Let $P\subseteq \pp^n$ be a fulldimensional convex polytope.
 Consider a linear space $L \subseteq \pp^n$ of dimension $l$ such that $\dim(L\cap P)< \dim(L)$. Then $L\cap P$ contains at most one face of $P$ of dimension $l-1$.
\end{lem}

\begin{thm} \label{thm:pointresidual}
 Let $P \subset \pp^n$ be a full-dimensional convex polytope with $d$ facets and let $f \in \C[X_0, \dots, X_n]_{d-n-1}$ be a polynomial satisfying
 \begin{equation*}
     \ord_x(f) \geq \ord_P(x)
 \end{equation*}
 for all $x \in \cR_0(P)$. Then
 \begin{equation*}
     \ord_x(f) \geq \ord_P(L_x)
 \end{equation*}
 for all $x \in \pp^n$.
\end{thm}

\begin{proof}[Proof of \Cref{thm:pointresidual}]
Let $f \in \C[X_0, \dots, X_n]_{d-n-1}$ vanish in the point-residual, i.e.
\begin{equation*} \label{eq:pointresidual}
    \mu_x(f) \geq \ord_P(x)
\end{equation*}
for all $x \in \cR_0(P)$. We want to prove that \begin{align*}
    \mu_{L_x}(f)\geq \ord_P(L_x)
\end{align*}must already hold for all $x \in \pp^n$. We will do so by induction on
\begin{equation*}
    r = \dim L_x = n-\rk(\cF_x).
\end{equation*}
The case $r = 0$ describes the vanishing in the point-residual, which is covered by our hypothesis. Thus let $r > 1$. Fix an enumeration $(\cF_i)_{i=1}^k$ of all flats of $M_P$ satisfying $\cF_x \subseteq \cF_i$ and $\rk(\cF_x) = \rk(\cF_i) - 1$. Geometrically, $L_i:=\bigcap\cF_i$ is a hyperplane in $L_x$.

We want to invoke \Cref{cor: van_order_Mult} to prove that $f$ vanishes in $x$ to the right order. As suggested by the definition of $\ord_P(L_x)$, we distinguish between two cases depending on whether the inequality $\dim(L_x\cap P)\leq \dim (L_x)$ is strict.
\begin{enumerate}
    \item Assume that $\dim(L_x\cap P)= \dim (L_x)$ and hence $\ord_P(L_x)=\nl_{M_P}(\cF_x)$. We wish to apply Corollary \ref{cor: van_order_Mult} to $D=d-n-1,m=\nl_{M_P}(\cF_x)$ and $H_i=L_i, i=1,\ldots,k$. We obtain that $\mu_{L_x}(f)\geq \nl_{M_P}(\cF_x)=\ord_P(L_x)$ is implied by 
    \begin{equation*}
    \begin{split}
        \sum_{i=1}^k \left(\mu_{L_i}(f) - (\nl_{M_P}(\cF_x) - 1)\right) &\geq d-n-1 - (\nl_{M_P}(\cF_x) - 1)
    \end{split}
    \end{equation*}
    To see that this is satisfied, we first note that by our inductive hypothesis and by definition of $\ord_P(L_i)$ we have
    \begin{equation*}
        \mu_{L_i}(f) \geq \ord_P(L_i) \geq \nl_{M_P}(\cF_i) 
    \end{equation*}
    for all $i$. Therefore
    \begin{equation*}
    \begin{split}
        \sum_i \left( \mu_{L_i}(f) - (\nl_{M_P}(\cF_x) - 1) \right) &\geq \sum_i \left( \nl_{M_P}(\cF_i) - \nl_{M_P}(\cF) + 1 \right) = \\
        &= \sum_i \left(|\cF_i| - |\cF_x| + \rk(\cF_x) - \rk(\cF_i) + 1\right) = \\
        &= \sum_i |\cF_i \setminus \cF_x| = d  - |\cF_x| > \\
        &> d - n - 1 - (\nl_{M_P}(\cF_x) - 1).
    \end{split}
    \end{equation*}
    The first equality in the penultimate line comes from the fact that the sets $\cF_i\setminus \cF_x$ form a partition of $\cH_P\setminus \cF_x$. For the last inequality we use that since $r > 0$ we have $n > \rk(\cF_x) = |\cF_x| - \nl_{M_P}(\cF_x)$.
    
    \item Now assume that $\dim(L_x\cap P)< \dim (L_x)$ and therefore $\ord_P(L_x)=\nl_{M_P}(\cF_x)+1$. This time we apply Corollary \ref{cor: van_order_Mult} to $D=d-n-1,m=\nl_{M_P}(\cF_x)+1$ and $H_i=L_i, i=1,\ldots,k$. We obtain that $\mu_{L_x}(f)\geq \nl_{M_P}+1(\cF_x)=\ord_P(L_x)$ is implied by 
    \begin{equation*}
    \begin{split}
        \sum_{i=1}^k \left(\mu_{L_i}(f) - \nl_{M_P}(\cF_x)\right) &\geq d-n-1 - \nl_{M_P}(\cF_x)
    \end{split}
    \end{equation*}

    As in case (1) we have
    \begin{equation*}
        \mu_{L_i}(f) \geq \ord_P(L_i) \geq \nl_{M_P}(\cF_i) 
    \end{equation*}
    for all $i$. The second inequality by definition is an equality precisely when $\dim(L_i\cap P)=\dim(L_i)$. By \Cref{lem: convpol_noninc} this can happen for at most one index $i\in \{1,\ldots,k\}$.  For the remaining indices, the inequality is strict. Therefore we conclude as in the first case:
    \begin{equation*}
    \begin{split}
        \sum_i \left( \mu_{L_i}(f) - \nl_{M_P}(\cF_x) \right) \geq &\sum_i \left( \nl_{M_P}(\cF_i) + 1 - \nl_{M_P}(\cF_x) \right) - 1 = \\
        &= \sum_i \left(|\cF_i| - |\cF_x|\right) - 1 = \sum_i \left( |\cF_i \setminus \cF_x| \right) - 1 = \\
        &= d - |\cF_x| - 1 \\&> d - n - 1 - \nl_{M_P}(\cF_x).
    \end{split} 
    \end{equation*}
\end{enumerate}
\end{proof}

\end{document}
